\documentclass[12pt]{article}

\usepackage{setspace}
\onehalfspace

%\usepackage[utf8]{inputenc}
%\usepackage{t1enc}
%\usepackage[T1]{fontenc}

\usepackage{amsmath}


%\DeclareMathOperator{\Exp}{Exp}
%\def\Exp{{\mathrm{Exp}}}

%\DeclareMathOperator{\Norm}{{\cal N}}
%\def\Norm{{\cal N}}

%\DeclareMathOperator{\Binom}{Binom}
%\def\Binom{{\mathrm{Binom}}}

%\DeclareMathOperator{\E}{E}
%\def\E{{\mathrm{E}}}

%\DeclareMathOperator{\D}{D}
%\def\D{{\mathrm{D}}}

%\renewcommand{\P}{\operatorname{P}}
%\def\P{{\mathrm{P}}}

%\DeclareMathOperator{\cov}{cov}
%\def\cov{{\mathrm{cov}}}


%\DeclareMathOperator{\corr}{corr}
%\def\corr{{\mathrm{corr}}}


%\renewcommand{\d}{\operatorname{d\! }} %integral \! a nyero
%\def\d{{\mathrm{d\! }}}


%\DeclareMathOperator{\R}{\mathbb{R}}
%\def\R{{\mathbb{R}}}

%\DeclareMathOperator{\Unif}{{\cal U}}
%\def\Unif{{\cal U}}

%\DeclareMathOperator{\Poi}{{Poisson}}
%\def\Poi{{\mathrm{Poisson}}}

%\def\ldots{{...\:}}


\usepackage{moodle}


\begin{document}

\begin{quiz}{gyakorló}

\begin{multi}[feedback={Például: legyen a kísérlet random húzás a $[0,1]$-ből és $A=\{\frac{1}{2}\}$}]{esem10}
Ha $\P(A)=0$ akkor $A$ sosem következik be.
\item igaz
\item* hamis
\end{multi}
\begin{multi}[feedback={A fenti formula a $u$-próba ún. próbastatisztikája.}]{proba8}
Az $t$-próbánál a 
$$
u=\frac{\overline{X}-\mu_0}{\sigma}\sqrt{n}.
$$
mennyiség segítségével döntünk a hipotéziseinkről.
\item igaz
\item* hamis
\end{multi}
\begin{multi}[feedback={$B=A\cup (B\setminus A)$+additívitás}]{val4}
A valószínűség monoton, azaz $\P(A)\le \P(B)$, ha $A\subseteq B$.
\item* igaz
\item hamis
\end{multi}
\begin{multi}[feedback={Definíció.}]{esem7}
Az $A$ és $B$ események pontosan akkor függetlenek, ha $\P(AB)=\P(A)\P(B)$.
\item* igaz
\item hamis
\end{multi}
\begin{multi}[feedback={Igen, éppen ez alapján tudunk dönteni. Ahhoz képest hogy sztenderd normális átlagos, tipikus értéket, vagy kiugró, extrém értéket látunk.}]{proba9}
Ha egy $16$ elemű minta normális sokaságból származik, $\sigma=3$ és hipotéziseink:
$$
H_0: \mu=10
$$
$$
H_1: \mu\neq 10
$$
akkor az 
$$
\frac{\overline{X}-10}{3}4
$$
valségi változó sztenderd normális eloszlású ha $H_0$ igaz.
\item* igaz
\item hamis
\end{multi}
\begin{multi}[feedback={Def.}]{felt2}
$A$-nak a $B$-re vonatkozó feltételes valségét $\P(A|B)=\frac{\P(A\cdot B)}{\P(B)}$ módon értelmezzük, feltéve hogy $\P(B)>0$.
\item* igaz
\item hamis
\end{multi}
\begin{multi}[feedback={Az $s^2$ torzítatlan becslése a sokasági (elméleti) szórásnégyzetnek}]{minta8}
Az $s^2$ korrigált empirikus szórásnégyzetre $\E(s^2)=\E(X)$.
\item igaz
\item* hamis
\end{multi}
\begin{multi}[feedback={Definíció}]{minta3}
Egy $X_1,\ldots,X_n\sim X$ minta esetén a mintaátlag $\overline{X}=\frac{X_1+\ldots+X_n}{n-1}$
\item igaz
\item* hamis
\end{multi}
\begin{multi}[feedback={Def.}]{sfv1}
Azt mondjuk, hogy egy $f:\R\to \R$ sűrűségfüggvény, ha 
$$
f\ge 0 \ \ \text{és}\ \ \int_{-\infty}^{+\infty}f(x)\d x=1
$$
\item* igaz
\item hamis
\end{multi}
\begin{multi}[feedback={Def.}]{impfvv10}
Ha $\xi\sim \Exp(\lambda)$ akkor a sűrűségfüggvénye 
$$
f(x)=
\begin{cases}
1-e^{-\lambda x}& \ \ x>0\\
0 & \text{máskor} 
\end{cases}
$$
\item igaz
\item* hamis
\end{multi}
\begin{multi}[feedback={Def.}]{impdvv6}
Azt mondjuk, hogy $\xi\sim \Binom(n,p)$, ha 
$$
\P(\xi=k)=\binom{n}{k}p^k(1-p)\ \ \ (k=0,1\ldots n)
$$
valamilyen $n>0$ és $p\in[0,1]$ esetén.
\item igaz
\item* hamis
\end{multi}
\begin{multi}[feedback={Páronként független valségi változókra a szórásnégyzet additív, általában nem.}]{dvv14}
A $\xi$ és $\eta$ valségi változókra:
$$
\D^2(\xi+\eta)=\D^2(\xi)+\D^2(\eta)
$$
feltéve, hogy létezik a szórásuk.
\item igaz
\item* hamis
\end{multi}
\begin{multi}[feedback={Tulajdonság.}]{flenvv6}
Ha $\xi$ és $\eta$ valségi változóknak létezik a szórása, akkor a 
kovarianciájuk
$$
\cov(\xi,\eta)=\E(\xi\eta)-\E(\xi)\E(\eta)
$$
\item* igaz
\item hamis
\end{multi}
\begin{multi}[feedback={Def.}]{dvv9}
Egy $\xi$ valségi változó az $x_1,\ldots, x_n$ értékeket veheti fel és $p_k=\P(\xi=x_k)$. 
Ekkor a második momentuma:
$$
\E(\xi^2)=\sum_{k=1}^n x_k p_k^2
$$
\item igaz
\item* hamis
\end{multi}
\begin{multi}[feedback={Fontos a $p_k\ge 0$ is.}]{dvv1}
Egy $p_1,\ldots, p_n$ vektor pont akkor valségeloszlás, ha $\sum_{k=1}^n p_k=1$.
\item igaz
\item* hamis
\end{multi}
\begin{multi}[feedback={Az $s^2$ torzítatlan becslése a sokasági (elméleti) szórásnégyzetnek}]{minta7}
Az $s^2$ korrigált empirikus szórásnégyzetre $\E(s^2)=\D^2(X)$.
\item* igaz
\item hamis
\end{multi}
\begin{multi}[feedback={Komplementer események.}]{eofv9}
Bármely $\xi$-re és $a$-ra $F_{\xi}(a)+\P(\xi\ge a)=1$.
\item* igaz
\item hamis
\end{multi}
\begin{multi}[feedback={Def.}]{eofv1}
Egy $F:\R\to \R$ pontosan akkor eloszlásfüggvény, ha monoton nemcsökkenő, 
balról folytonos, $\lim_{x\to+\infty}F(x)=1$ és $\lim_{x\to-\infty}F(x)=0$
\item* igaz
\item hamis
\end{multi}
\begin{multi}[feedback={$\xi=\xi_1+\ldots+\xi_n$, ahol $\xi_k$ Bernoulli és $\E(\xi_k)=p$+$\E$-additívitása}]{impdvv3}
Egy $\xi\sim \Binom(n,p)$ szórásnégyzete: $np(1-p)$
\item* igaz
\item hamis
\end{multi}
\begin{multi}[feedback={A szórást ismerni kell.}]{proba2}
Ha a minta normális sokaságból származik és a sokasági várható értékre (átlagra) vonatkozó kérdésről 
akarunk dönteni akkor $u$-próbát használhatunk.
\item igaz
\item* hamis
\end{multi}
\begin{multi}[feedback={Def.+$P(AB)=P(A)$ a feltételekből.}]{felt1}
$A$-nak a $B$-re ($\P(B)>0$) vonatkozó feltételes valsége $\P(A|B)=\frac{\P(A)}{\P(B)}$, ha $A$ maga után vonja $B$-t.
\item* igaz
\item hamis
\end{multi}
\begin{multi}[feedback={Def.}]{sfv8}
Ha van $\xi$-nek $f$ sűrűségfüggvénye, akkor bármely $b$-re:
$$
\P(\xi<b)=f(b)
$$
\item igaz
\item* hamis
\end{multi}
\begin{multi}[feedback={Ezek az $u$-próba használatának előzményei.}]{proba1}
Ha a minta normális sokaságból származik, ismert a szórás és a sokasági várható értékre (átlagra) vonatkozó kérdésről 
akarunk dönteni akkor $u$-próbát használhatunk.
\item* igaz
\item hamis
\end{multi}
\begin{multi}[feedback={A fenti formula az $u$-próba ún. próbastatisztikája.}]{proba5}
Az $u$-próbánál a 
$$
u=\frac{\overline{X}-\mu_0}{\sigma}\sqrt{n}
$$
mennyiség segítségével döntünk a hipotéziseinkről.
\item* igaz
\item hamis
\end{multi}
\begin{multi}[feedback={Additívitás.}]{val6}
Bármely esemény valsége kiszámolható a $\frac{\text{kedvező}}{\text{összes}}$ képlettel, ha az elemi események 
egyforma valségűek.
\item* igaz
\item hamis
\end{multi}
\begin{multi}[feedback={$\xi=\xi_1+\ldots+\xi_n$, ahol $\xi_k$ Bernoulli és $\E(\xi)=p$+additívitás}]{impdvv1}
Egy $\xi\sim \Binom(n,p)$ várható értéke: $np$
\item* igaz
\item hamis
\end{multi}
\begin{multi}[feedback={$f\ge 0$ fontos (monoton nemcsökkenő $F$-nek a deriváltja)}]{sfv3}
Azt mondjuk, hogy egy $f:\R\to \R$ sűrűségfüggvény 
$$
\int_{-\infty}^{+\infty}f(x)\d x\ge 0
$$
\item igaz
\item* hamis
\end{multi}
\begin{multi}[feedback={Def.}]{felt3}
$A_1,\ldots,A_n$ teljes eseményrendszer, ha $\sum A_k=\Omega$ és a tagok páronként kizáróak.
\item* igaz
\item hamis
\end{multi}
\begin{multi}[feedback={Nókoment.}]{esem5}
Ha $\P(A)=1$ akkor $A$ bekövetkezhet.
\item* igaz
\item hamis
\end{multi}
\begin{multi}[feedback={Def.}]{sfv11}
Ha $\xi$-nek $F$ illetve $f$ az eloszlásfüggvény illetve sűrűségfüggvénye, akkor bármely $a<b$-re:
$$
F(b)-F(a)=\int_{a}^{b}f(x)\d x
$$
\item* igaz
\item hamis
\end{multi}
\begin{multi}[feedback={wiki}]{impfvv3}
Egy $\xi\sim \Unif(a,b)$ szórásnégyzete: $\frac{(b-a)^2}{12}$
\item* igaz
\item hamis
\end{multi}
\begin{multi}[feedback={de-Morgan azonosság}]{esem2}
Az $\overline{A+B}$ és $\overline{A} \cdot \overline{B}$ események kizárják egymást.
\item igaz
\item* hamis
\end{multi}
\begin{multi}[feedback={Egy konstans v.v-nak 0 a szórása.}]{dvv12}
Egy $\xi$ valségi változó szórása mindig pozitív.
\item igaz
\item* hamis
\end{multi}
\begin{multi}[feedback={Def.}]{impfvv8}
Ha $\xi\sim \Norm(0,1)$ akkor a sűrűségfüggvénye $f(x)=\frac{e^{-\frac{x^2}{2}}}{\sqrt{2\pi}}$.
\item* igaz
\item hamis
\end{multi}
\begin{multi}[feedback={Def.}]{flenvv5}
Ha $\xi$ és $\eta$ diszkrét valségi változóknak létezik a szórása és nemnulla, akkor a 
korrelációjuk
$$
\corr(\xi,\eta)=\E(\xi\eta)-\E(\xi)\E(\eta)
$$
\item igaz
\item* hamis
\end{multi}
\begin{multi}[feedback={Def.}]{impfvv11}
Ha $\xi\sim \Unif(a,b)$ akkor a sűrűségfüggvénye 
$$
f(x)=
\begin{cases}
\frac{1}{b-a} & \ \ a<x<b \\
0 & \text{máskor} 
\end{cases}
$$
\item* igaz
\item hamis
\end{multi}
\begin{multi}[feedback={Def.}]{dvv8}
Egy $\xi$ valségi változó az $x_1,\ldots, x_n$ értékeket veheti fel és $p_k=\P(\xi=x_k)$. 
Ekkor a második momentuma:
$$
\E(\xi^2)=\sum_{k=1}^n x_k^2 p_k
$$
\item* igaz
\item hamis
\end{multi}
\begin{multi}[feedback={Def.}]{impfvv9}
Ha $\xi\sim \Exp(\lambda)$ akkor a sűrűségfüggvénye 
$$
f(x)=
\begin{cases}
\lambda e^{-\lambda x}& \ \ x>0\\
0 & \text{máskor} 
\end{cases}
$$
\item* igaz
\item hamis
\end{multi}
\begin{multi}[feedback={Def.}]{flenvv4}
Ha $\xi$ és $\eta$ diszkrét valségi változóknak létezik a szórása és nemnulla, akkor a 
korrelációjuk
$$
\corr(\xi,\eta)=\frac{\cov(\xi,\eta)}{\D(\xi)\D(\eta)}
$$
\item* igaz
\item hamis
\end{multi}
\begin{multi}[feedback={Csak ha az elemi események egyforma valségűek.}]{val5}
Bármely esemény valsége kiszámolható a $\frac{\text{kedvező}}{\text{összes}}$ képlettel.
\item igaz
\item* hamis
\end{multi}
\begin{multi}[feedback={Def./tulajdonság}]{dvv10}
Egy $\xi$ valségi változó az $x_1,\ldots, x_n$ értékeket veheti fel. 
Ekkor a szórásnégyzete:
$$
\D^2(\xi)=\E(\xi^2)-\E(\xi)^2
$$
\item* igaz
\item hamis
\end{multi}
\begin{multi}[feedback={Az $F(x)=\P(\xi<x)$ konstans az $(a,b)$-n ha $a$ és $b$ két szomszédos eleme az értékkészletnek.}]{eofv6}
Véges értékkészletű valségi változók eloszlásfüggvénye lépcsős függvény.
\item* igaz
\item hamis
\end{multi}
\begin{multi}[feedback={Mindegyiknek van, definíció szerint.}]{eofv5}
A diszkrét valségi változóknak nincsen eloszlásfüggvénye.
\item igaz
\item* hamis
\end{multi}
\begin{multi}[feedback={Tulajdonság,tétel.}]{flenvv9}
Ha a $\xi$ és $\eta$ valségi változóknak létezik a korrelációja, akkor $-1\le \corr(\xi,\eta)\le 1$.
\item* igaz
\item hamis
\end{multi}
\end{quiz}

\end{document}

